\begin{table*}[t]\centering\small\caption{Oracle ablation: temperature fitted on in-distribution validation data (TS) versus on validation data shifted like the test condition (Orc). The ID-fitted $T$ is constant across shifts by construction.}\label{tab:oracle}
\begin{tabular}{lcccccc}
\toprule
Shift & ECE TS & ECE Orc & NLL TS & NLL Orc & $T$ TS & $T$ Orc \\
\midrule
r0 & \rhval{main/mlp-+-temperature-scaling/rot0/ece/mean:3} & \rhval{abl_oracle/ts-oracle-shifted-val/rot0/ece/mean:3} & \rhval{main/mlp-+-temperature-scaling/rot0/nll/mean:3} & \rhval{abl_oracle/ts-oracle-shifted-val/rot0/nll/mean:3} & \rhval{main/mlp-+-temperature-scaling/rot0/temperature/mean:3} & \rhval{abl_oracle/ts-oracle-shifted-val/rot0/temperature/mean:3} \\
r10 & \rhval{main/mlp-+-temperature-scaling/rot10/ece/mean:3} & \rhval{abl_oracle/ts-oracle-shifted-val/rot10/ece/mean:3} & \rhval{main/mlp-+-temperature-scaling/rot10/nll/mean:3} & \rhval{abl_oracle/ts-oracle-shifted-val/rot10/nll/mean:3} & \rhval{main/mlp-+-temperature-scaling/rot10/temperature/mean:3} & \rhval{abl_oracle/ts-oracle-shifted-val/rot10/temperature/mean:3} \\
r20 & \rhval{main/mlp-+-temperature-scaling/rot20/ece/mean:3} & \rhval{abl_oracle/ts-oracle-shifted-val/rot20/ece/mean:3} & \rhval{main/mlp-+-temperature-scaling/rot20/nll/mean:3} & \rhval{abl_oracle/ts-oracle-shifted-val/rot20/nll/mean:3} & \rhval{main/mlp-+-temperature-scaling/rot20/temperature/mean:3} & \rhval{abl_oracle/ts-oracle-shifted-val/rot20/temperature/mean:3} \\
r30 & \rhval{main/mlp-+-temperature-scaling/rot30/ece/mean:3} & \rhval{abl_oracle/ts-oracle-shifted-val/rot30/ece/mean:3} & \rhval{main/mlp-+-temperature-scaling/rot30/nll/mean:3} & \rhval{abl_oracle/ts-oracle-shifted-val/rot30/nll/mean:3} & \rhval{main/mlp-+-temperature-scaling/rot30/temperature/mean:3} & \rhval{abl_oracle/ts-oracle-shifted-val/rot30/temperature/mean:3} \\
r45 & \rhval{main/mlp-+-temperature-scaling/rot45/ece/mean:3} & \rhval{abl_oracle/ts-oracle-shifted-val/rot45/ece/mean:3} & \rhval{main/mlp-+-temperature-scaling/rot45/nll/mean:3} & \rhval{abl_oracle/ts-oracle-shifted-val/rot45/nll/mean:3} & \rhval{main/mlp-+-temperature-scaling/rot45/temperature/mean:3} & \rhval{abl_oracle/ts-oracle-shifted-val/rot45/temperature/mean:3} \\
r60 & \rhval{main/mlp-+-temperature-scaling/rot60/ece/mean:3} & \rhval{abl_oracle/ts-oracle-shifted-val/rot60/ece/mean:3} & \rhval{main/mlp-+-temperature-scaling/rot60/nll/mean:3} & \rhval{abl_oracle/ts-oracle-shifted-val/rot60/nll/mean:3} & \rhval{main/mlp-+-temperature-scaling/rot60/temperature/mean:3} & \rhval{abl_oracle/ts-oracle-shifted-val/rot60/temperature/mean:3} \\
n0.25 & \rhval{main/mlp-+-temperature-scaling/noise0.25/ece/mean:3} & \rhval{abl_oracle/ts-oracle-shifted-val/noise0.25/ece/mean:3} & \rhval{main/mlp-+-temperature-scaling/noise0.25/nll/mean:3} & \rhval{abl_oracle/ts-oracle-shifted-val/noise0.25/nll/mean:3} & \rhval{main/mlp-+-temperature-scaling/noise0.25/temperature/mean:3} & \rhval{abl_oracle/ts-oracle-shifted-val/noise0.25/temperature/mean:3} \\
n0.5 & \rhval{main/mlp-+-temperature-scaling/noise0.5/ece/mean:3} & \rhval{abl_oracle/ts-oracle-shifted-val/noise0.5/ece/mean:3} & \rhval{main/mlp-+-temperature-scaling/noise0.5/nll/mean:3} & \rhval{abl_oracle/ts-oracle-shifted-val/noise0.5/nll/mean:3} & \rhval{main/mlp-+-temperature-scaling/noise0.5/temperature/mean:3} & \rhval{abl_oracle/ts-oracle-shifted-val/noise0.5/temperature/mean:3} \\
n0.75 & \rhval{main/mlp-+-temperature-scaling/noise0.75/ece/mean:3} & \rhval{abl_oracle/ts-oracle-shifted-val/noise0.75/ece/mean:3} & \rhval{main/mlp-+-temperature-scaling/noise0.75/nll/mean:3} & \rhval{abl_oracle/ts-oracle-shifted-val/noise0.75/nll/mean:3} & \rhval{main/mlp-+-temperature-scaling/noise0.75/temperature/mean:3} & \rhval{abl_oracle/ts-oracle-shifted-val/noise0.75/temperature/mean:3} \\
n1.0 & \rhval{main/mlp-+-temperature-scaling/noise1.0/ece/mean:3} & \rhval{abl_oracle/ts-oracle-shifted-val/noise1.0/ece/mean:3} & \rhval{main/mlp-+-temperature-scaling/noise1.0/nll/mean:3} & \rhval{abl_oracle/ts-oracle-shifted-val/noise1.0/nll/mean:3} & \rhval{main/mlp-+-temperature-scaling/noise1.0/temperature/mean:3} & \rhval{abl_oracle/ts-oracle-shifted-val/noise1.0/temperature/mean:3} \\
r30+n0.5 & \rhval{main/mlp-+-temperature-scaling/rot30_noise0.5/ece/mean:3} & \rhval{abl_oracle/ts-oracle-shifted-val/rot30_noise0.5/ece/mean:3} & \rhval{main/mlp-+-temperature-scaling/rot30_noise0.5/nll/mean:3} & \rhval{abl_oracle/ts-oracle-shifted-val/rot30_noise0.5/nll/mean:3} & \rhval{main/mlp-+-temperature-scaling/rot30_noise0.5/temperature/mean:3} & \rhval{abl_oracle/ts-oracle-shifted-val/rot30_noise0.5/temperature/mean:3} \\
\bottomrule
\end{tabular}
\end{table*}
\begin{table*}[t]\centering\small\caption{Dropout decomposition (exploratory): MLP, the dropout-trained network evaluated deterministically with dropout off (Det), and MC dropout with 20 passes (MCD). Means over 5 seeds.}\label{tab:dropdet}
\begin{tabular}{lcccccccc}
\toprule
Shift & Acc MLP & Acc Det & Acc MCD & ECE MLP & ECE Det & ECE MCD & NLL Det & NLL MCD \\
\midrule
r0 & \rhval{main/mlp/rot0/accuracy/mean:3} & \rhval{abl_dropdet/dropout-trained-mlp-deterministic/rot0/accuracy/mean:3} & \rhval{main/mc-dropout/rot0/accuracy/mean:3} & \rhval{main/mlp/rot0/ece/mean:3} & \rhval{abl_dropdet/dropout-trained-mlp-deterministic/rot0/ece/mean:3} & \rhval{main/mc-dropout/rot0/ece/mean:3} & \rhval{abl_dropdet/dropout-trained-mlp-deterministic/rot0/nll/mean:3} & \rhval{main/mc-dropout/rot0/nll/mean:3} \\
r10 & \rhval{main/mlp/rot10/accuracy/mean:3} & \rhval{abl_dropdet/dropout-trained-mlp-deterministic/rot10/accuracy/mean:3} & \rhval{main/mc-dropout/rot10/accuracy/mean:3} & \rhval{main/mlp/rot10/ece/mean:3} & \rhval{abl_dropdet/dropout-trained-mlp-deterministic/rot10/ece/mean:3} & \rhval{main/mc-dropout/rot10/ece/mean:3} & \rhval{abl_dropdet/dropout-trained-mlp-deterministic/rot10/nll/mean:3} & \rhval{main/mc-dropout/rot10/nll/mean:3} \\
r20 & \rhval{main/mlp/rot20/accuracy/mean:3} & \rhval{abl_dropdet/dropout-trained-mlp-deterministic/rot20/accuracy/mean:3} & \rhval{main/mc-dropout/rot20/accuracy/mean:3} & \rhval{main/mlp/rot20/ece/mean:3} & \rhval{abl_dropdet/dropout-trained-mlp-deterministic/rot20/ece/mean:3} & \rhval{main/mc-dropout/rot20/ece/mean:3} & \rhval{abl_dropdet/dropout-trained-mlp-deterministic/rot20/nll/mean:3} & \rhval{main/mc-dropout/rot20/nll/mean:3} \\
r30 & \rhval{main/mlp/rot30/accuracy/mean:3} & \rhval{abl_dropdet/dropout-trained-mlp-deterministic/rot30/accuracy/mean:3} & \rhval{main/mc-dropout/rot30/accuracy/mean:3} & \rhval{main/mlp/rot30/ece/mean:3} & \rhval{abl_dropdet/dropout-trained-mlp-deterministic/rot30/ece/mean:3} & \rhval{main/mc-dropout/rot30/ece/mean:3} & \rhval{abl_dropdet/dropout-trained-mlp-deterministic/rot30/nll/mean:3} & \rhval{main/mc-dropout/rot30/nll/mean:3} \\
r45 & \rhval{main/mlp/rot45/accuracy/mean:3} & \rhval{abl_dropdet/dropout-trained-mlp-deterministic/rot45/accuracy/mean:3} & \rhval{main/mc-dropout/rot45/accuracy/mean:3} & \rhval{main/mlp/rot45/ece/mean:3} & \rhval{abl_dropdet/dropout-trained-mlp-deterministic/rot45/ece/mean:3} & \rhval{main/mc-dropout/rot45/ece/mean:3} & \rhval{abl_dropdet/dropout-trained-mlp-deterministic/rot45/nll/mean:3} & \rhval{main/mc-dropout/rot45/nll/mean:3} \\
r60 & \rhval{main/mlp/rot60/accuracy/mean:3} & \rhval{abl_dropdet/dropout-trained-mlp-deterministic/rot60/accuracy/mean:3} & \rhval{main/mc-dropout/rot60/accuracy/mean:3} & \rhval{main/mlp/rot60/ece/mean:3} & \rhval{abl_dropdet/dropout-trained-mlp-deterministic/rot60/ece/mean:3} & \rhval{main/mc-dropout/rot60/ece/mean:3} & \rhval{abl_dropdet/dropout-trained-mlp-deterministic/rot60/nll/mean:3} & \rhval{main/mc-dropout/rot60/nll/mean:3} \\
n0.25 & \rhval{main/mlp/noise0.25/accuracy/mean:3} & \rhval{abl_dropdet/dropout-trained-mlp-deterministic/noise0.25/accuracy/mean:3} & \rhval{main/mc-dropout/noise0.25/accuracy/mean:3} & \rhval{main/mlp/noise0.25/ece/mean:3} & \rhval{abl_dropdet/dropout-trained-mlp-deterministic/noise0.25/ece/mean:3} & \rhval{main/mc-dropout/noise0.25/ece/mean:3} & \rhval{abl_dropdet/dropout-trained-mlp-deterministic/noise0.25/nll/mean:3} & \rhval{main/mc-dropout/noise0.25/nll/mean:3} \\
n0.5 & \rhval{main/mlp/noise0.5/accuracy/mean:3} & \rhval{abl_dropdet/dropout-trained-mlp-deterministic/noise0.5/accuracy/mean:3} & \rhval{main/mc-dropout/noise0.5/accuracy/mean:3} & \rhval{main/mlp/noise0.5/ece/mean:3} & \rhval{abl_dropdet/dropout-trained-mlp-deterministic/noise0.5/ece/mean:3} & \rhval{main/mc-dropout/noise0.5/ece/mean:3} & \rhval{abl_dropdet/dropout-trained-mlp-deterministic/noise0.5/nll/mean:3} & \rhval{main/mc-dropout/noise0.5/nll/mean:3} \\
n0.75 & \rhval{main/mlp/noise0.75/accuracy/mean:3} & \rhval{abl_dropdet/dropout-trained-mlp-deterministic/noise0.75/accuracy/mean:3} & \rhval{main/mc-dropout/noise0.75/accuracy/mean:3} & \rhval{main/mlp/noise0.75/ece/mean:3} & \rhval{abl_dropdet/dropout-trained-mlp-deterministic/noise0.75/ece/mean:3} & \rhval{main/mc-dropout/noise0.75/ece/mean:3} & \rhval{abl_dropdet/dropout-trained-mlp-deterministic/noise0.75/nll/mean:3} & \rhval{main/mc-dropout/noise0.75/nll/mean:3} \\
n1.0 & \rhval{main/mlp/noise1.0/accuracy/mean:3} & \rhval{abl_dropdet/dropout-trained-mlp-deterministic/noise1.0/accuracy/mean:3} & \rhval{main/mc-dropout/noise1.0/accuracy/mean:3} & \rhval{main/mlp/noise1.0/ece/mean:3} & \rhval{abl_dropdet/dropout-trained-mlp-deterministic/noise1.0/ece/mean:3} & \rhval{main/mc-dropout/noise1.0/ece/mean:3} & \rhval{abl_dropdet/dropout-trained-mlp-deterministic/noise1.0/nll/mean:3} & \rhval{main/mc-dropout/noise1.0/nll/mean:3} \\
r30+n0.5 & \rhval{main/mlp/rot30_noise0.5/accuracy/mean:3} & \rhval{abl_dropdet/dropout-trained-mlp-deterministic/rot30_noise0.5/accuracy/mean:3} & \rhval{main/mc-dropout/rot30_noise0.5/accuracy/mean:3} & \rhval{main/mlp/rot30_noise0.5/ece/mean:3} & \rhval{abl_dropdet/dropout-trained-mlp-deterministic/rot30_noise0.5/ece/mean:3} & \rhval{main/mc-dropout/rot30_noise0.5/ece/mean:3} & \rhval{abl_dropdet/dropout-trained-mlp-deterministic/rot30_noise0.5/nll/mean:3} & \rhval{main/mc-dropout/rot30_noise0.5/nll/mean:3} \\
\bottomrule
\end{tabular}
\end{table*}
\begin{table}[t]\centering\small\caption{Sensitivity sweeps (mean over seeds 0--2, conditions r0, r30, n0.5): ensemble size $M$ and dropout rate $p$. The $M=5$ and $p=\rhval{const/dropout_p:1}$ rows are deliberate re-runs of the main configuration on 3 of the 5 seeds.}\label{tab:sweeps}
\resizebox{\columnwidth}{!}{\begin{tabular}{llcccccc}
\toprule
Sweep & Value & ECE r0 & NLL r0 & ECE r30 & NLL r30 & ECE n0.5 & NLL n0.5 \\
\midrule
M & 1 & \rhval{sweep_members/deep-ensemble@members=1/rot0/ece/mean:3} & \rhval{sweep_members/deep-ensemble@members=1/rot0/nll/mean:3} & \rhval{sweep_members/deep-ensemble@members=1/rot30/ece/mean:3} & \rhval{sweep_members/deep-ensemble@members=1/rot30/nll/mean:3} & \rhval{sweep_members/deep-ensemble@members=1/noise0.5/ece/mean:3} & \rhval{sweep_members/deep-ensemble@members=1/noise0.5/nll/mean:3} \\
M & 2 & \rhval{sweep_members/deep-ensemble@members=2/rot0/ece/mean:3} & \rhval{sweep_members/deep-ensemble@members=2/rot0/nll/mean:3} & \rhval{sweep_members/deep-ensemble@members=2/rot30/ece/mean:3} & \rhval{sweep_members/deep-ensemble@members=2/rot30/nll/mean:3} & \rhval{sweep_members/deep-ensemble@members=2/noise0.5/ece/mean:3} & \rhval{sweep_members/deep-ensemble@members=2/noise0.5/nll/mean:3} \\
M & 3 & \rhval{sweep_members/deep-ensemble@members=3/rot0/ece/mean:3} & \rhval{sweep_members/deep-ensemble@members=3/rot0/nll/mean:3} & \rhval{sweep_members/deep-ensemble@members=3/rot30/ece/mean:3} & \rhval{sweep_members/deep-ensemble@members=3/rot30/nll/mean:3} & \rhval{sweep_members/deep-ensemble@members=3/noise0.5/ece/mean:3} & \rhval{sweep_members/deep-ensemble@members=3/noise0.5/nll/mean:3} \\
M & 5 & \rhval{sweep_members/deep-ensemble@members=5/rot0/ece/mean:3} & \rhval{sweep_members/deep-ensemble@members=5/rot0/nll/mean:3} & \rhval{sweep_members/deep-ensemble@members=5/rot30/ece/mean:3} & \rhval{sweep_members/deep-ensemble@members=5/rot30/nll/mean:3} & \rhval{sweep_members/deep-ensemble@members=5/noise0.5/ece/mean:3} & \rhval{sweep_members/deep-ensemble@members=5/noise0.5/nll/mean:3} \\
M & 10 & \rhval{sweep_members/deep-ensemble@members=10/rot0/ece/mean:3} & \rhval{sweep_members/deep-ensemble@members=10/rot0/nll/mean:3} & \rhval{sweep_members/deep-ensemble@members=10/rot30/ece/mean:3} & \rhval{sweep_members/deep-ensemble@members=10/rot30/nll/mean:3} & \rhval{sweep_members/deep-ensemble@members=10/noise0.5/ece/mean:3} & \rhval{sweep_members/deep-ensemble@members=10/noise0.5/nll/mean:3} \\
\midrule
$p$ & 0.1 & \rhval{sweep_dropout/mc-dropout@p=0.1/rot0/ece/mean:3} & \rhval{sweep_dropout/mc-dropout@p=0.1/rot0/nll/mean:3} & \rhval{sweep_dropout/mc-dropout@p=0.1/rot30/ece/mean:3} & \rhval{sweep_dropout/mc-dropout@p=0.1/rot30/nll/mean:3} & \rhval{sweep_dropout/mc-dropout@p=0.1/noise0.5/ece/mean:3} & \rhval{sweep_dropout/mc-dropout@p=0.1/noise0.5/nll/mean:3} \\
$p$ & 0.2 & \rhval{sweep_dropout/mc-dropout@p=0.2/rot0/ece/mean:3} & \rhval{sweep_dropout/mc-dropout@p=0.2/rot0/nll/mean:3} & \rhval{sweep_dropout/mc-dropout@p=0.2/rot30/ece/mean:3} & \rhval{sweep_dropout/mc-dropout@p=0.2/rot30/nll/mean:3} & \rhval{sweep_dropout/mc-dropout@p=0.2/noise0.5/ece/mean:3} & \rhval{sweep_dropout/mc-dropout@p=0.2/noise0.5/nll/mean:3} \\
$p$ & 0.3 & \rhval{sweep_dropout/mc-dropout@p=0.3/rot0/ece/mean:3} & \rhval{sweep_dropout/mc-dropout@p=0.3/rot0/nll/mean:3} & \rhval{sweep_dropout/mc-dropout@p=0.3/rot30/ece/mean:3} & \rhval{sweep_dropout/mc-dropout@p=0.3/rot30/nll/mean:3} & \rhval{sweep_dropout/mc-dropout@p=0.3/noise0.5/ece/mean:3} & \rhval{sweep_dropout/mc-dropout@p=0.3/noise0.5/nll/mean:3} \\
$p$ & 0.5 & \rhval{sweep_dropout/mc-dropout@p=0.5/rot0/ece/mean:3} & \rhval{sweep_dropout/mc-dropout@p=0.5/rot0/nll/mean:3} & \rhval{sweep_dropout/mc-dropout@p=0.5/rot30/ece/mean:3} & \rhval{sweep_dropout/mc-dropout@p=0.5/rot30/nll/mean:3} & \rhval{sweep_dropout/mc-dropout@p=0.5/noise0.5/ece/mean:3} & \rhval{sweep_dropout/mc-dropout@p=0.5/noise0.5/nll/mean:3} \\
\bottomrule
\end{tabular}
}\end{table}
\begin{figure}[t]\centering\includegraphics[width=\columnwidth]{sweeps.pdf}\caption{Sweeps over ensemble size (top) and dropout rate (bottom); mean $\pm$ SD over seeds 0--2.}\label{fig:sweeps}\end{figure}

\textbf{Oracle temperature (tests H5): supported, with a caveat.} Fitting $T$ on shift-matched validation data lowers ECE relative to the ID-fitted $T$ at every severe condition tested (r45, r60, n0.75, n1.0), e.g.\ r60 \rhval{abl_oracle/ts-oracle-shifted-val/rot60/ece/mean:3} versus \rhval{main/mlp-+-temperature-scaling/rot60/ece/mean:3} (the largest paired $p$ of the four is \rhval{cmp/abl_oracle/ts-oracle-shifted-val/noise0.75/ece/paired_p:sci1}, at n0.75, Table~\ref{tab:hyp}), and over all eleven conditions the oracle ECE stays at or below \rhval{abl_oracle/ts-oracle-shifted-val/rot30/ece/mean:3} (Table~\ref{tab:oracle}). The fitted oracle temperature grows with intensity, from \rhval{abl_oracle/ts-oracle-shifted-val/rot0/temperature/mean:3} at r0 to \rhval{abl_oracle/ts-oracle-shifted-val/rot60/temperature/mean:3} at r60 (Fig.~\ref{fig:conf}, Table~\ref{tab:oracle}), whereas the ID-fitted $T$ stays at \rhval{main/mlp-+-temperature-scaling/rot0/temperature/mean:3}. So the miscalibration of ID-fitted TS under shift is a temperature mismatch rather than a limit of scalar rescaling \emph{for ECE}. The caveat is in the NLL columns: the oracle NLL at r60 is \rhval{abl_oracle/ts-oracle-shifted-val/rot60/nll/mean:3}, close to the $\ln 10$ of a uniform prediction over ten classes, i.e.\ the oracle ``calibrates'' by making predictions nearly uniform. That is calibrated but uninformative; accuracy is unchanged. The oracle also needs to know the shift, which is not available in deployment, and the very large temperatures are fitted on only \rhval{const/n_val} validation images. We cannot say from this study whether a practical shift-agnostic procedure, such as perturbing the calibration set~\cite{tomani2020post}, would find them.

\textbf{Why does MC dropout win? A decomposition (exploratory).} The dropout-trained network evaluated deterministically (Det) is not identical to the MLP: its ECE is between MLP and MCD at most conditions (e.g.\ r30 \rhval{main/mlp/rot30/ece/mean:3}, \rhval{abl_dropdet/dropout-trained-mlp-deterministic/rot30/ece/mean:3}, \rhval{main/mc-dropout/rot30/ece/mean:3}; n1.0 \rhval{main/mlp/noise1.0/ece/mean:3}, \rhval{abl_dropdet/dropout-trained-mlp-deterministic/noise1.0/ece/mean:3}, \rhval{main/mc-dropout/noise1.0/ece/mean:3}; Table~\ref{tab:dropdet}), while at r0 the three are close (\rhval{main/mlp/rot0/ece/mean:3}, \rhval{abl_dropdet/dropout-trained-mlp-deterministic/rot0/ece/mean:3} and \rhval{main/mc-dropout/rot0/ece/mean:3}). Accuracy of Det is higher than the MLP at r10 (\rhval{abl_dropdet/dropout-trained-mlp-deterministic/rot10/accuracy/mean:3} vs \rhval{main/mlp/rot10/accuracy/mean:3}) and n0.25 (\rhval{abl_dropdet/dropout-trained-mlp-deterministic/noise0.25/accuracy/mean:3} vs \rhval{main/mlp/noise0.25/accuracy/mean:3}). The paired Det-versus-MLP ECE differences are not significant at r0 ($p=\rhval{cmp/abl_dropdet/mlp/rot0/ece/paired_p:2}$), r30 ($p=\rhval{cmp/abl_dropdet/mlp/rot30/ece/paired_p:3}$) or r60 ($p=\rhval{cmp/abl_dropdet/mlp/rot60/ece/paired_p:2}$) and only at n1.0 ($p=\rhval{cmp/abl_dropdet/mlp/noise1.0/ece/paired_p:3}$), whereas Det versus MCD is significant at r30, r60 and n1.0 (the largest $p$ is \rhval{cmp/abl_dropdet/mc-dropout/rot60/ece/paired_p:3}, at r60) but not at r0 ($p=\rhval{cmp/abl_dropdet/mc-dropout/rot0/ece/paired_p:2}$). So the evidence that MC averaging at test time contributes is clearer than that for dropout regularisation, which may also contribute, and a part of MC dropout's advantage over the ensemble is a property of the differently trained network rather than of its test-time Bayesian approximation (cf.~\cite{folgoc2021is}). We did not test a dropout-regularised ensemble, which would be the natural controlled comparison.

\textbf{Sensitivity to ensemble size and dropout rate.} In Table~\ref{tab:sweeps} the ensemble's NLL under r30 falls monotonically from \rhval{sweep_members/deep-ensemble@members=1/rot30/nll/mean:3} ($M=1$) to \rhval{sweep_members/deep-ensemble@members=10/rot30/nll/mean:3} ($M=10$), and under n0.5 from \rhval{sweep_members/deep-ensemble@members=1/noise0.5/nll/mean:3} to \rhval{sweep_members/deep-ensemble@members=10/noise0.5/nll/mean:3}; ECE falls from \rhval{sweep_members/deep-ensemble@members=1/rot30/ece/mean:3} to \rhval{sweep_members/deep-ensemble@members=10/rot30/ece/mean:3} (r30) and \rhval{sweep_members/deep-ensemble@members=1/noise0.5/ece/mean:3} to \rhval{sweep_members/deep-ensemble@members=10/noise0.5/ece/mean:3} (n0.5). Even $M=10$ stays above MC dropout at $p=\rhval{const/dropout_p:1}$ on r30 (NLL \rhval{sweep_dropout/mc-dropout@p=0.2/rot30/nll/mean:3}, ECE \rhval{sweep_dropout/mc-dropout@p=0.2/rot30/ece/mean:3}) and n0.5 (NLL \rhval{sweep_dropout/mc-dropout@p=0.2/noise0.5/nll/mean:3}, ECE \rhval{sweep_dropout/mc-dropout@p=0.2/noise0.5/ece/mean:3}), so the H3/H4 verdicts are not an artefact of $M=5$. Larger dropout rates reduce shifted ECE and NLL further ($p=0.5$: r30 ECE \rhval{sweep_dropout/mc-dropout@p=0.5/rot30/ece/mean:3}, n0.5 ECE \rhval{sweep_dropout/mc-dropout@p=0.5/noise0.5/ece/mean:3}) but worsen in-distribution calibration (r0 ECE \rhval{sweep_dropout/mc-dropout@p=0.5/rot0/ece/mean:3} versus \rhval{sweep_dropout/mc-dropout@p=0.2/rot0/ece/mean:3} at $p=\rhval{const/dropout_p:1}$), a trade-off that the main comparison at $p=\rhval{const/dropout_p:1}$ does not reveal. The sweeps use seeds 0--2 only; they are not used to choose $p$ or $M$, and a rate chosen on shifted data would be tuning on the shift.
